%% LyX 2.4.4 created this file.  For more info, see https://www.lyx.org/.
%% Do not edit unless you really know what you are doing.
\documentclass[english]{article}
\usepackage[T1]{fontenc}
\usepackage[latin9]{inputenc}
\usepackage{enumitem}
\usepackage{amsmath}
\usepackage{amsthm}
\usepackage{geometry}
\geometry{verbose,tmargin=1in,bmargin=1in,lmargin=1in,rmargin=1in}

\makeatletter
%%%%%%%%%%%%%%%%%%%%%%%%%%%%%% Textclass specific LaTeX commands.
\numberwithin{equation}{section}
\numberwithin{figure}{section}
\newlength{\lyxlabelwidth}      % auxiliary length 

%%%%%%%%%%%%%%%%%%%%%%%%%%%%%% User specified LaTeX commands.
\usepackage{fancyhdr}
\usepackage{lastpage}
\pagestyle{fancy}
\usepackage{enumitem}
\usepackage{ifthen}
%sets math fonts to be more readable
\everymath=\expandafter{\the\everymath\displaystyle}
%\DeclareMathSizes{10}{10}{10}{10}
\usepackage{multicol}

\usepackage{tikz}
\usetikzlibrary{patterns}
\usetikzlibrary{plotmarks}
\usepackage{pgfplots}
\pgfplotsset{compat=newest}

\usepackage{advdate}    % Advancing/saving dates
\usepackage{datetime}   % Dates formatting
\usepackage{datenumber} % Counters for dates
\usdate
% Added by lyx2lyx
\setlength{\parskip}{\smallskipamount}
\setlength{\parindent}{0pt}

\makeatother

\usepackage{babel}
\begin{document}
\renewcommand{\author}{Dr.\,James Ashe}

\newcommand{\classNum}{137}

\newcommand{\classSec}{B}

\newcommand{\nameType}{}

\newcommand{\examType}{Exam 2}

\renewcommand{\date}{will be be given on \formatdate{13}{3}{2013}} %%\currenttime, \today

%%First page's header%%%%%%%%%%%%%%%%%%%%%%
\fancypagestyle{firststyle} %%header settings for first pagr
{
	\fancyhead[L]{MTH \classNum{}(\classSec{}) \author{}\\
	\textbf{\examType{}} \date}
	\fancyhead[C]{\textbf{\nameType}}
	\setlength{\headsep}{14pt}
}
%%%%%%%%%%%%%%%%%%%%%%%%%%%%%%%%%%%%%%%%%%

%%Next page's header%%%%%%%%%%%%%%%%%%%%%%%
\setlist{nolistsep}
\lhead{\classNum{}(\classSec{})} 
\chead{\textbf{\nameType}} 
\cfoot{\thepage{}~of~\pageref{LastPage}} 
\setlength{\headsep}{5pt} 
%%%%%%%%%%%%%%%%%%%%%%%%%%%%%%%%%%%%%%%%%%%

%%Various settings; see the preamble for other settings%%%%%
%%\setenumerate[0]{leftmargin=12pt,itemindent=0pt,labelwidth=0pt}
\renewcommand{\labelenumi}{\textbf{\arabic{enumi}.}}
\renewcommand{\labelenumii}{\textbf{\alph{enumii}.}}
\thispagestyle{firststyle}
%%%%%%%%%%%%%%%%%%%%%%%%%%%%%%%%%%%%%%%%%%

\vfill$ $

\subsection*{2.1 Functions}
\begin{itemize}
\item Determine if a given relation is a function.

\begin{itemize}
\item From a graph; \textbf{HW 3-4.}
\item From a list of coordinates; \textbf{HW 5-6, 9-12}
\item From an equation; \textbf{HW 7-8}
\end{itemize}
\item Determine the domain and range of a relation.

\begin{itemize}
\item From a graph. \emph{see the quiz}
\item From a list; \textbf{HW 9-10}
\end{itemize}
\item Understanding and working with function notation, i.e., $f(x)$.

\begin{itemize}
\item From a list; \textbf{HW 11-12}
\item From an equation; \textbf{HW 13-14}
\end{itemize}
\item Find the average rate of change; \textbf{HW 15-16}
\item Find the difference quotient; \textbf{HW 17-19}
\end{itemize}

\subsection*{2.2: Graphing functions}
\begin{itemize}
\item Understand how to read and write interval notation, e.g., $1\leq x<5$
means $x$ is in $[1,5)$.
\item Sketch a relation from an equation, then give its domain and range.

\begin{itemize}
\item Common functions like lines, quadratics, and absolute values; \textbf{HW
1-3, 5}
\item Piecewise functions; \textbf{HW 6-8, 15-17}
\end{itemize}
\item Identify when a function is decreasing, increasing, or constant for
all type of functions including piecewise functions; \textbf{HW 10-14}
\end{itemize}

\subsection*{2.3: Function transformations.}
\begin{itemize}
\item Understand and graph functions using the parent function and basic
transformations. \textbf{HW 1-20}

\begin{itemize}
\item Parent functions include: $f(x)=x$, $f(x)=x^{2}$, $f(x)=\sqrt{x}$,
and $f(x)=\left|x\right|$.
\item Transformations:

\begin{itemize}
\item $f(x\pm h)$: $"+"\Rightarrow Left$ and $"-"\Rightarrow Right$
\item $f(x)\pm k$: $"+"\Rightarrow Up$ and $"-"\Rightarrow Down$
\item $a\cdot f(x)$: $a>0\Rightarrow normal$ and $a<0\Rightarrow"flip\mbox{ }upside\mbox{ }down"$
\item $a\cdot f(x)$: $\left|a\right|>1\Rightarrow Narrow$ and $\left|a\right|<1\Rightarrow Wide$
\end{itemize}
\end{itemize}
\end{itemize}

\subsection*{2.4 Function operations.}
\begin{itemize}
\item Add, subtract, multiply, and divide functions; \textbf{HW 1-8}

\begin{itemize}
\item This includes finding the general function and the value at specific
points, e.g., $f-g$ and $(f-g)(2)$.
\item You may also need to find the domain of the new function. 
\end{itemize}
\item Function compositions; $(f\circ g)(x)=f(g(x))$; \textbf{HW 8-10,
12-13, 16-19}
\end{itemize}

\subsection*{2.5 Inverse function}
\begin{itemize}
\item Determine if a function is $1-1$/invertible, from a list or graph;
\textbf{HW 1-9}
\item Find the inverse function $f^{-1}$; \textbf{HW 10, 12, 16-20}
\item Sketch an inverse function from a graph, or determine if two graphs
are inverses of each other; \textbf{HW 13-16}
\item Understand inverse function notation.\vfill
\end{itemize}

\end{document}
